%%%%%%%%%%%%%%%%%%%%%%%%%%%%%%%%%%%%%%%%%%%%%%%%%%%%%%%%%%%%
% 8.6 HYPOTHESIS TESTING
% The Composition of Advanced Mathematics
%%%%%%%%%%%%%%%%%%%%%%%%%%%%%%%%%%%%%%%%%%%%%%%%%%%%%%%%%%%%

\section{Hypothesis Testing}
\label{sec:hypothesis-testing}

\prerequisite{
Sampling distributions, statistical estimation, confidence intervals,
normal distributions, Student's \(t\)-distribution, chi-square and
\(F\)-distributions, likelihood, standard errors, and basic probability.
}

\learningobjectives{
By the end of this section, the reader should be able to:
\begin{itemize}
    \item define null and alternative hypotheses;
    \item distinguish simple and composite hypotheses;
    \item distinguish one-sided and two-sided tests;
    \item define a test statistic;
    \item define a rejection region;
    \item define a significance level;
    \item explain Type I and Type II errors;
    \item define test power;
    \item interpret \(p\)-values correctly;
    \item distinguish \(p\)-values from posterior probabilities;
    \item perform one-sample \(z\)-tests for means;
    \item perform one-sample \(t\)-tests for means;
    \item perform tests for one population proportion;
    \item perform two-sample tests for independent means;
    \item distinguish pooled and Welch two-sample tests;
    \item perform paired \(t\)-tests;
    \item perform tests for differences of proportions;
    \item perform tests for normal variances;
    \item perform tests for ratios of normal variances;
    \item understand chi-square goodness-of-fit tests;
    \item understand chi-square tests of independence;
    \item understand expected cell counts;
    \item understand likelihood-ratio tests;
    \item understand score and Wald tests conceptually;
    \item understand the Neyman--Pearson framework;
    \item understand uniformly most powerful tests conceptually;
    \item compute and interpret power;
    \item understand sample-size planning for tests;
    \item distinguish statistical significance from practical
          significance;
    \item understand multiple-testing problems;
    \item apply Bonferroni corrections;
    \item distinguish family-wise error rate and false discovery rate;
    \item understand permutation and randomization tests conceptually;
    \item connect hypothesis testing with confidence intervals,
          estimation, regression, ANOVA, and experimental design.
\end{itemize}
}

%%%%%%%%%%%%%%%%%%%%%%%%%%%%%%%%%%%%%%%%%%%%%%%%%%%%%%%%%%%%
\subsection{What Is a Statistical Hypothesis?}
\label{subsec:what-is-statistical-hypothesis}
%%%%%%%%%%%%%%%%%%%%%%%%%%%%%%%%%%%%%%%%%%%%%%%%%%%%%%%%%%%%

A statistical hypothesis is a statement about a population parameter or
probability model.

Examples include

\[
\mu=50,
\]

\[
p=0.40,
\]

or

\[
\sigma_1^2=\sigma_2^2.
\]

Hypothesis testing provides a formal procedure for evaluating how
compatible observed data are with such a claim.

%%%%%%%%%%%%%%%%%%%%%%%%%%%%%%%%%%%%%%%%%%%%%%%%%%%%%%%%%%%%
\subsection{Null Hypothesis}
\label{subsec:null-hypothesis}
%%%%%%%%%%%%%%%%%%%%%%%%%%%%%%%%%%%%%%%%%%%%%%%%%%%%%%%%%%%%

\begin{definitionbox}{Null Hypothesis}
The null hypothesis, denoted \(H_0\), is the reference claim tested
against the data.
\end{definitionbox}

Examples include

\[
\boxed{
H_0:\mu=\mu_0
}
\]

or

\[
\boxed{
H_0:p=p_0.
}
\]

The null hypothesis determines the sampling distribution used to measure
how unusual the observed statistic is.

%%%%%%%%%%%%%%%%%%%%%%%%%%%%%%%%%%%%%%%%%%%%%%%%%%%%%%%%%%%%
\subsection{Alternative Hypothesis}
\label{subsec:alternative-hypothesis}
%%%%%%%%%%%%%%%%%%%%%%%%%%%%%%%%%%%%%%%%%%%%%%%%%%%%%%%%%%%%

\begin{definitionbox}{Alternative Hypothesis}
The alternative hypothesis, denoted \(H_A\) or \(H_1\), describes the
parameter values regarded as competing with the null hypothesis.
\end{definitionbox}

Common forms include

\[
\boxed{
H_A:\mu\neq\mu_0,
}
\]

\[
\boxed{
H_A:\mu>\mu_0,
}
\]

or

\[
\boxed{
H_A:\mu<\mu_0.
}
\]

%%%%%%%%%%%%%%%%%%%%%%%%%%%%%%%%%%%%%%%%%%%%%%%%%%%%%%%%%%%%
\subsection{Simple and Composite Hypotheses}
\label{subsec:simple-composite-hypotheses}
%%%%%%%%%%%%%%%%%%%%%%%%%%%%%%%%%%%%%%%%%%%%%%%%%%%%%%%%%%%%

A simple hypothesis specifies the parameter completely.

For example,

\[
\boxed{
H_0:\mu=10.
}
\]

A composite hypothesis includes several possible parameter values.

For example,

\[
\boxed{
H_A:\mu>10.
}
\]

Thus a one-point claim is simple, while a range of parameter values is
composite.

%%%%%%%%%%%%%%%%%%%%%%%%%%%%%%%%%%%%%%%%%%%%%%%%%%%%%%%%%%%%
\subsection{Two-Sided Tests}
\label{subsec:two-sided-tests}
%%%%%%%%%%%%%%%%%%%%%%%%%%%%%%%%%%%%%%%%%%%%%%%%%%%%%%%%%%%%

A two-sided alternative has the form

\[
\boxed{
H_A:\theta\neq\theta_0.
}
\]

Evidence against \(H_0\) may appear in either direction.

The rejection region therefore lies in both tails of the null sampling
distribution.

%%%%%%%%%%%%%%%%%%%%%%%%%%%%%%%%%%%%%%%%%%%%%%%%%%%%%%%%%%%%
\subsection{Right-Tailed Tests}
\label{subsec:right-tailed-tests}
%%%%%%%%%%%%%%%%%%%%%%%%%%%%%%%%%%%%%%%%%%%%%%%%%%%%%%%%%%%%

A right-tailed alternative has the form

\[
\boxed{
H_A:\theta>\theta_0.
}
\]

Large values of the test statistic provide evidence against the null.

The rejection region lies in the upper tail.

%%%%%%%%%%%%%%%%%%%%%%%%%%%%%%%%%%%%%%%%%%%%%%%%%%%%%%%%%%%%
\subsection{Left-Tailed Tests}
\label{subsec:left-tailed-tests}
%%%%%%%%%%%%%%%%%%%%%%%%%%%%%%%%%%%%%%%%%%%%%%%%%%%%%%%%%%%%

A left-tailed alternative has the form

\[
\boxed{
H_A:\theta<\theta_0.
}
\]

Small values of the test statistic provide evidence against the null.

The rejection region lies in the lower tail.

%%%%%%%%%%%%%%%%%%%%%%%%%%%%%%%%%%%%%%%%%%%%%%%%%%%%%%%%%%%%
\subsection{The Test Statistic}
\label{subsec:test-statistic}
%%%%%%%%%%%%%%%%%%%%%%%%%%%%%%%%%%%%%%%%%%%%%%%%%%%%%%%%%%%%

A test statistic is a numerical summary of the data chosen so its
distribution under \(H_0\) is known or approximately known.

A common standardized form is

\[
\boxed{
T
=
\frac{
\text{estimate}
-
\text{null value}
}{
\text{standard error under }H_0
}.
}
\]

Large departures from zero may indicate inconsistency with the null.

%%%%%%%%%%%%%%%%%%%%%%%%%%%%%%%%%%%%%%%%%%%%%%%%%%%%%%%%%%%%
\subsection{Rejection Region}
\label{subsec:rejection-region}
%%%%%%%%%%%%%%%%%%%%%%%%%%%%%%%%%%%%%%%%%%%%%%%%%%%%%%%%%%%%

\begin{definitionbox}{Rejection Region}
The rejection region is the set of test-statistic values for which the
null hypothesis is rejected.
\end{definitionbox}

The rejection region is chosen before observing the test outcome.

Its probability under the null is controlled by the significance level.

%%%%%%%%%%%%%%%%%%%%%%%%%%%%%%%%%%%%%%%%%%%%%%%%%%%%%%%%%%%%
\subsection{Significance Level}
\label{subsec:significance-level}
%%%%%%%%%%%%%%%%%%%%%%%%%%%%%%%%%%%%%%%%%%%%%%%%%%%%%%%%%%%%

The significance level is denoted

\[
\boxed{
\alpha.
}
\]

It is the maximum or specified probability of rejecting the null when
the null is true under the testing procedure.

Common choices are

\[
\boxed{
0.10,
\quad
0.05,
\quad
0.01.
}
\]

%%%%%%%%%%%%%%%%%%%%%%%%%%%%%%%%%%%%%%%%%%%%%%%%%%%%%%%%%%%%
\subsection{Type I Error}
\label{subsec:type-i-error}
%%%%%%%%%%%%%%%%%%%%%%%%%%%%%%%%%%%%%%%%%%%%%%%%%%%%%%%%%%%%

\begin{definitionbox}{Type I Error}
A Type I error occurs when \(H_0\) is rejected even though \(H_0\) is
true.
\end{definitionbox}

Symbolically,

\[
\boxed{
\text{Type I Error}
=
\text{reject }H_0
\text{ when }H_0\text{ is true}.
}
\]

For a level-\(\alpha\) test,

\[
\boxed{
P(\text{Type I error})
\leq
\alpha,
}
\]

with equality in many standard continuous testing procedures.

%%%%%%%%%%%%%%%%%%%%%%%%%%%%%%%%%%%%%%%%%%%%%%%%%%%%%%%%%%%%
\subsection{Type II Error}
\label{subsec:type-ii-error}
%%%%%%%%%%%%%%%%%%%%%%%%%%%%%%%%%%%%%%%%%%%%%%%%%%%%%%%%%%%%

\begin{definitionbox}{Type II Error}
A Type II error occurs when the test fails to reject \(H_0\) even though
a specified alternative is true.
\end{definitionbox}

Its probability is often denoted

\[
\boxed{
\beta.
}
\]

Unlike \(\alpha\), \(\beta\) generally depends on the particular
alternative parameter value.

%%%%%%%%%%%%%%%%%%%%%%%%%%%%%%%%%%%%%%%%%%%%%%%%%%%%%%%%%%%%
\subsection{Power}
\label{subsec:test-power}
%%%%%%%%%%%%%%%%%%%%%%%%%%%%%%%%%%%%%%%%%%%%%%%%%%%%%%%%%%%%

\begin{definitionbox}{Power}
The power of a test at parameter value \(\theta\) is the probability of
rejecting \(H_0\) when that parameter value is true.
\end{definitionbox}

Thus,

\[
\boxed{
\operatorname{Power}(\theta)
=
P_\theta(\text{reject }H_0).
}
\]

At a specific alternative,

\[
\boxed{
\operatorname{Power}
=
1-\beta.
}
\]

High power means the test is likely to detect meaningful departures from
the null.

%%%%%%%%%%%%%%%%%%%%%%%%%%%%%%%%%%%%%%%%%%%%%%%%%%%%%%%%%%%%
\subsection{Decision Table}
\label{subsec:hypothesis-decision-table}
%%%%%%%%%%%%%%%%%%%%%%%%%%%%%%%%%%%%%%%%%%%%%%%%%%%%%%%%%%%%

The basic outcomes are

\[
\begin{array}{c|cc}
&
H_0\text{ true}
&
H_0\text{ false}
\\
\hline
\text{Reject }H_0
&
\text{Type I error}
&
\text{Correct detection}
\\
\text{Fail to reject }H_0
&
\text{Correct decision}
&
\text{Type II error}
\end{array}
\]

Statistical tests balance the risks of these errors.

%%%%%%%%%%%%%%%%%%%%%%%%%%%%%%%%%%%%%%%%%%%%%%%%%%%%%%%%%%%%
\subsection{Fail to Reject Versus Accept}
\label{subsec:fail-to-reject}
%%%%%%%%%%%%%%%%%%%%%%%%%%%%%%%%%%%%%%%%%%%%%%%%%%%%%%%%%%%%

A nonsignificant result is usually described as

\[
\boxed{
\text{fail to reject }H_0
}
\]

rather than

\[
\boxed{
\text{accept }H_0.
}
\]

Failing to find strong evidence against a hypothesis is not the same as
proving the hypothesis true.

A study may simply have insufficient precision or power.

%%%%%%%%%%%%%%%%%%%%%%%%%%%%%%%%%%%%%%%%%%%%%%%%%%%%%%%%%%%%
\subsection{\(p\)-Value}
\label{subsec:p-value}
%%%%%%%%%%%%%%%%%%%%%%%%%%%%%%%%%%%%%%%%%%%%%%%%%%%%%%%%%%%%

\begin{definitionbox}{\(p\)-Value}
The \(p\)-value is the probability, computed under the null hypothesis,
of obtaining a test statistic at least as extreme as the observed value
in the direction specified by the alternative.
\end{definitionbox}

Thus,

\[
\boxed{
p\text{-value}
=
P_{H_0}
(
\text{result at least as incompatible with }H_0
).
}
\]

%%%%%%%%%%%%%%%%%%%%%%%%%%%%%%%%%%%%%%%%%%%%%%%%%%%%%%%%%%%%
\subsection{Interpretation of a Small \(p\)-Value}
\label{subsec:small-p-value}
%%%%%%%%%%%%%%%%%%%%%%%%%%%%%%%%%%%%%%%%%%%%%%%%%%%%%%%%%%%%

A small \(p\)-value means the observed data would be relatively unusual
under the null model.

Thus the data provide evidence against \(H_0\).

A conventional decision rule is

\[
\boxed{
p\leq\alpha
\Longrightarrow
\text{reject }H_0.
}
\]

If

\[
p>\alpha,
\]

fail to reject \(H_0\).

%%%%%%%%%%%%%%%%%%%%%%%%%%%%%%%%%%%%%%%%%%%%%%%%%%%%%%%%%%%%
\subsection{What a \(p\)-Value Is Not}
\label{subsec:p-value-not}
%%%%%%%%%%%%%%%%%%%%%%%%%%%%%%%%%%%%%%%%%%%%%%%%%%%%%%%%%%%%

A \(p\)-value is not

\[
\boxed{
P(H_0\text{ is true}\mid\text{data}).
}
\]

It is also not

\[
\boxed{
P(\text{results occurred by chance}).
}
\]

The \(p\)-value is calculated under the assumption that the null model
is true.

Posterior probabilities of hypotheses require a Bayesian model.

%%%%%%%%%%%%%%%%%%%%%%%%%%%%%%%%%%%%%%%%%%%%%%%%%%%%%%%%%%%%
\subsection{Statistical Significance}
\label{subsec:statistical-significance}
%%%%%%%%%%%%%%%%%%%%%%%%%%%%%%%%%%%%%%%%%%%%%%%%%%%%%%%%%%%%

A result is called statistically significant at level \(\alpha\) if

\[
\boxed{
p\leq\alpha.
}
\]

This means the result crosses a predetermined evidence threshold.

It does not automatically imply that the effect is:

\[
\boxed{
\text{large},
}
\]

\[
\boxed{
\text{important},
}
\]

or

\[
\boxed{
\text{causal}.
}
\]

%%%%%%%%%%%%%%%%%%%%%%%%%%%%%%%%%%%%%%%%%%%%%%%%%%%%%%%%%%%%
\subsection{Practical Significance}
\label{subsec:practical-significance}
%%%%%%%%%%%%%%%%%%%%%%%%%%%%%%%%%%%%%%%%%%%%%%%%%%%%%%%%%%%%

Practical significance concerns whether the estimated effect is large
enough to matter scientifically, economically, medically, or
operationally.

A very large sample can make a tiny effect statistically significant.

Thus results should be interpreted using both

\[
\boxed{
\text{effect size}
}
\]

and

\[
\boxed{
\text{uncertainty}.
}
\]

Confidence intervals are especially useful for this purpose.

%%%%%%%%%%%%%%%%%%%%%%%%%%%%%%%%%%%%%%%%%%%%%%%%%%%%%%%%%%%%
\subsection{Testing a Population Mean with Known \(\sigma\)}
\label{subsec:z-test-one-mean}
%%%%%%%%%%%%%%%%%%%%%%%%%%%%%%%%%%%%%%%%%%%%%%%%%%%%%%%%%%%%

Suppose

\[
X_1,\ldots,X_n
\]

come from a normal population with known \(\sigma\).

Test

\[
H_0:\mu=\mu_0.
\]

Under \(H_0\),

\[
\boxed{
Z
=
\frac{
\overline{X}-\mu_0
}{
\sigma/\sqrt{n}
}
\sim
N(0,1).
}
\]

This is the one-sample \(z\)-test.

%%%%%%%%%%%%%%%%%%%%%%%%%%%%%%%%%%%%%%%%%%%%%%%%%%%%%%%%%%%%
\subsection{Two-Sided \(z\)-Test}
\label{subsec:two-sided-z-test}
%%%%%%%%%%%%%%%%%%%%%%%%%%%%%%%%%%%%%%%%%%%%%%%%%%%%%%%%%%%%

For

\[
H_A:\mu\neq\mu_0,
\]

reject at level \(\alpha\) when

\[
\boxed{
|Z|
>
z_{\alpha/2},
}
\]

using upper-tail critical-value notation.

Equivalently, reject when the two-sided \(p\)-value satisfies

\[
\boxed{
2P(Z_{\text{std}}\geq|z_{\text{obs}}|)
\leq
\alpha.
}
\]

%%%%%%%%%%%%%%%%%%%%%%%%%%%%%%%%%%%%%%%%%%%%%%%%%%%%%%%%%%%%
\subsection{Worked Example: One-Sample \(z\)-Test}
\label{subsec:worked-one-sample-z-test}
%%%%%%%%%%%%%%%%%%%%%%%%%%%%%%%%%%%%%%%%%%%%%%%%%%%%%%%%%%%%

\begin{workedexamplebox}{Testing a Population Mean}

Suppose

\[
H_0:\mu=100,
\]

\[
H_A:\mu\neq100,
\]

with

\[
\sigma=15,
\qquad
n=100,
\qquad
\overline{x}=104.
\]

The test statistic is

\[
z
=
\frac{
104-100
}{
15/\sqrt{100}
}.
\]

Therefore,

\[
z
=
\frac4{1.5}
\approx
2.667.
\]

\begin{answerbox}
\[
\boxed{
z\approx2.667.
}
\]

The two-sided \(p\)-value is

\[
\boxed{
2[1-\Phi(2.667)].
}
\]
\end{answerbox}

At \(\alpha=0.05\), this result is significant because
\(|z|>1.96\).

\end{workedexamplebox}

%%%%%%%%%%%%%%%%%%%%%%%%%%%%%%%%%%%%%%%%%%%%%%%%%%%%%%%%%%%%
\subsection{Testing a Mean with Unknown \(\sigma\)}
\label{subsec:one-sample-t-test}
%%%%%%%%%%%%%%%%%%%%%%%%%%%%%%%%%%%%%%%%%%%%%%%%%%%%%%%%%%%%

Suppose

\[
X_i
\overset{\text{i.i.d.}}{\sim}
N(\mu,\sigma^2)
\]

with unknown \(\sigma\).

To test

\[
H_0:\mu=\mu_0,
\]

use

\[
\boxed{
T
=
\frac{
\overline{X}-\mu_0
}{
S/\sqrt{n}
}.
}
\]

Under \(H_0\),

\[
\boxed{
T
\sim
t_{n-1}.
}
\]

%%%%%%%%%%%%%%%%%%%%%%%%%%%%%%%%%%%%%%%%%%%%%%%%%%%%%%%%%%%%
\subsection{Worked Example: One-Sample \(t\)-Test}
\label{subsec:worked-one-sample-t-test}
%%%%%%%%%%%%%%%%%%%%%%%%%%%%%%%%%%%%%%%%%%%%%%%%%%%%%%%%%%%%

\begin{workedexamplebox}{Testing an Unknown-Variance Mean}

Suppose

\[
n=25,
\qquad
\overline{x}=54,
\qquad
s=10.
\]

Test

\[
H_0:\mu=50
\]

against

\[
H_A:\mu>50.
\]

Then

\[
t
=
\frac{
54-50
}{
10/\sqrt{25}
}
=
\frac4{2}
=
2.
\]

The degrees of freedom are

\[
24.
\]

\begin{answerbox}
\[
\boxed{
t=2,
\qquad
df=24.
}
\]

The right-tailed \(p\)-value is

\[
\boxed{
P(T_{24}\geq2).
}
\]
\end{answerbox}

\end{workedexamplebox}

%%%%%%%%%%%%%%%%%%%%%%%%%%%%%%%%%%%%%%%%%%%%%%%%%%%%%%%%%%%%
\subsection{Conditions for the One-Sample \(t\)-Test}
\label{subsec:t-test-conditions}
%%%%%%%%%%%%%%%%%%%%%%%%%%%%%%%%%%%%%%%%%%%%%%%%%%%%%%%%%%%%

The exact one-sample \(t\)-test assumes normal sampling.

For larger samples, the procedure is often approximately valid under
weaker distributional assumptions.

Severe outliers or heavy tails can still cause problems.

The sampling design must also justify the assumed independence structure.

%%%%%%%%%%%%%%%%%%%%%%%%%%%%%%%%%%%%%%%%%%%%%%%%%%%%%%%%%%%%
\subsection{Testing a Population Proportion}
\label{subsec:one-proportion-test}
%%%%%%%%%%%%%%%%%%%%%%%%%%%%%%%%%%%%%%%%%%%%%%%%%%%%%%%%%%%%

Suppose

\[
X\sim\operatorname{Bin}(n,p).
\]

To test

\[
H_0:p=p_0,
\]

the large-sample score-type test statistic is

\[
\boxed{
Z
=
\frac{
\widehat{p}-p_0
}{
\sqrt{
p_0(1-p_0)/n
}
}.
}
\]

Notice that the denominator uses the null value \(p_0\), not
\(\widehat{p}\).

%%%%%%%%%%%%%%%%%%%%%%%%%%%%%%%%%%%%%%%%%%%%%%%%%%%%%%%%%%%%
\subsection{Worked Example: One-Proportion Test}
\label{subsec:worked-one-proportion-test}
%%%%%%%%%%%%%%%%%%%%%%%%%%%%%%%%%%%%%%%%%%%%%%%%%%%%%%%%%%%%

\begin{workedexamplebox}{Testing a Success Probability}

Suppose

\[
n=200,
\qquad
x=130.
\]

Then

\[
\widehat{p}
=
0.65.
\]

Test

\[
H_0:p=0.60
\]

against

\[
H_A:p>0.60.
\]

The test statistic is

\[
z
=
\frac{
0.65-0.60
}{
\sqrt{
(0.60)(0.40)/200
}
}.
\]

\begin{answerbox}
\[
\boxed{
z
=
\frac{0.05}{
\sqrt{0.24/200}
}.
}
\]
\end{answerbox}

The right-tailed \(p\)-value is

\[
1-\Phi(z).
\]

\end{workedexamplebox}

%%%%%%%%%%%%%%%%%%%%%%%%%%%%%%%%%%%%%%%%%%%%%%%%%%%%%%%%%%%%
\subsection{Exact Binomial Test}
\label{subsec:exact-binomial-test}
%%%%%%%%%%%%%%%%%%%%%%%%%%%%%%%%%%%%%%%%%%%%%%%%%%%%%%%%%%%%

For small samples or proportions near \(0\) or \(1\), one may use the
exact binomial distribution rather than a normal approximation.

For a right-tailed test,

\[
\boxed{
p\text{-value}
=
P_{p_0}(X\geq x_{\text{obs}}).
}
\]

For a left-tailed test,

\[
\boxed{
p\text{-value}
=
P_{p_0}(X\leq x_{\text{obs}}).
}
\]

Two-sided exact \(p\)-values require a specified definition of
``at least as extreme.''

%%%%%%%%%%%%%%%%%%%%%%%%%%%%%%%%%%%%%%%%%%%%%%%%%%%%%%%%%%%%
\subsection{Two Independent Means}
\label{subsec:test-two-independent-means}
%%%%%%%%%%%%%%%%%%%%%%%%%%%%%%%%%%%%%%%%%%%%%%%%%%%%%%%%%%%%

Suppose two independent populations have means

\[
\mu_1
\]

and

\[
\mu_2.
\]

A common null hypothesis is

\[
\boxed{
H_0:
\mu_1-\mu_2
=
\Delta_0.
}
\]

Usually,

\[
\Delta_0=0.
\]

The point estimate is

\[
\overline{X}_1-\overline{X}_2.
\]

%%%%%%%%%%%%%%%%%%%%%%%%%%%%%%%%%%%%%%%%%%%%%%%%%%%%%%%%%%%%
\subsection{Welch Two-Sample \(t\)-Test}
\label{subsec:welch-two-sample-test}
%%%%%%%%%%%%%%%%%%%%%%%%%%%%%%%%%%%%%%%%%%%%%%%%%%%%%%%%%%%%

Without assuming equal population variances, use

\[
\boxed{
T
=
\frac{
(\overline{X}_1-\overline{X}_2)-\Delta_0
}{
\sqrt{
S_1^2/n_1
+
S_2^2/n_2
}
}.
}
\]

The approximate degrees of freedom are

\[
\boxed{
\nu
=
\frac{
\left(
S_1^2/n_1
+
S_2^2/n_2
\right)^2
}{
\frac{
(S_1^2/n_1)^2
}{
n_1-1
}
+
\frac{
(S_2^2/n_2)^2
}{
n_2-1
}
}.
}
\]

%%%%%%%%%%%%%%%%%%%%%%%%%%%%%%%%%%%%%%%%%%%%%%%%%%%%%%%%%%%%
\subsection{Worked Example: Welch Test}
\label{subsec:worked-welch-test}
%%%%%%%%%%%%%%%%%%%%%%%%%%%%%%%%%%%%%%%%%%%%%%%%%%%%%%%%%%%%

\begin{workedexamplebox}{Comparing Two Means}

Suppose

\[
\overline{x}_1=82,
\quad
s_1=12,
\quad
n_1=50,
\]

and

\[
\overline{x}_2=76,
\quad
s_2=10,
\quad
n_2=45.
\]

Test

\[
H_0:\mu_1-\mu_2=0.
\]

The statistic is

\[
t
=
\frac{
82-76
}{
\sqrt{
12^2/50
+
10^2/45
}
}.
\]

\begin{answerbox}
\[
\boxed{
t
=
\frac{
6
}{
\sqrt{
144/50+100/45
}
}.
}
\]
\end{answerbox}

The \(p\)-value is obtained from a \(t\)-distribution using Welch's
degrees of freedom.

\end{workedexamplebox}

%%%%%%%%%%%%%%%%%%%%%%%%%%%%%%%%%%%%%%%%%%%%%%%%%%%%%%%%%%%%
\subsection{Pooled Two-Sample \(t\)-Test}
\label{subsec:pooled-two-sample-test}
%%%%%%%%%%%%%%%%%%%%%%%%%%%%%%%%%%%%%%%%%%%%%%%%%%%%%%%%%%%%

If the two populations are assumed normal with common variance

\[
\sigma_1^2=\sigma_2^2,
\]

define

\[
\boxed{
S_p^2
=
\frac{
(n_1-1)S_1^2
+
(n_2-1)S_2^2
}{
n_1+n_2-2
}.
}
\]

Then

\[
\boxed{
T
=
\frac{
(\overline{X}_1-\overline{X}_2)-\Delta_0
}{
S_p
\sqrt{
1/n_1+1/n_2
}
}
}
\]

has

\[
\boxed{
n_1+n_2-2
}
\]

degrees of freedom under the null.

%%%%%%%%%%%%%%%%%%%%%%%%%%%%%%%%%%%%%%%%%%%%%%%%%%%%%%%%%%%%
\subsection{Welch Versus Pooled Testing}
\label{subsec:welch-versus-pooled-test}
%%%%%%%%%%%%%%%%%%%%%%%%%%%%%%%%%%%%%%%%%%%%%%%%%%%%%%%%%%%%

Welch's test does not assume equal population variances.

The pooled test does.

Therefore Welch's procedure is often preferred as a general default.

The pooled test can be efficient when the common-variance model is
appropriate, but its extra assumption should be justified rather than
made automatically.

%%%%%%%%%%%%%%%%%%%%%%%%%%%%%%%%%%%%%%%%%%%%%%%%%%%%%%%%%%%%
\subsection{Paired \(t\)-Test}
\label{subsec:paired-t-test}
%%%%%%%%%%%%%%%%%%%%%%%%%%%%%%%%%%%%%%%%%%%%%%%%%%%%%%%%%%%%

For paired data, define

\[
D_i
=
X_i-Y_i.
\]

Then test

\[
\boxed{
H_0:\mu_D=\Delta_0.
}
\]

The statistic is

\[
\boxed{
T
=
\frac{
\overline{D}-\Delta_0
}{
S_D/\sqrt{n}
}.
}
\]

Under normal differences,

\[
\boxed{
T\sim t_{n-1}.
}
\]

%%%%%%%%%%%%%%%%%%%%%%%%%%%%%%%%%%%%%%%%%%%%%%%%%%%%%%%%%%%%
\subsection{Worked Example: Paired \(t\)-Test}
\label{subsec:worked-paired-t-test}
%%%%%%%%%%%%%%%%%%%%%%%%%%%%%%%%%%%%%%%%%%%%%%%%%%%%%%%%%%%%

\begin{workedexamplebox}{Before-and-After Difference}

Suppose

\[
n=16,
\qquad
\overline{d}=3,
\qquad
s_D=4.
\]

Test

\[
H_0:\mu_D=0
\]

against

\[
H_A:\mu_D>0.
\]

Then

\[
t
=
\frac{
3
}{
4/\sqrt{16}
}
=
3.
\]

\begin{answerbox}
\[
\boxed{
t=3,
\qquad
df=15.
}
\]
\end{answerbox}

The \(p\)-value is

\[
P(T_{15}\geq3).
\]

\end{workedexamplebox}

%%%%%%%%%%%%%%%%%%%%%%%%%%%%%%%%%%%%%%%%%%%%%%%%%%%%%%%%%%%%
\subsection{Testing Difference of Two Proportions}
\label{subsec:test-two-proportions}
%%%%%%%%%%%%%%%%%%%%%%%%%%%%%%%%%%%%%%%%%%%%%%%%%%%%%%%%%%%%

Suppose independent samples produce

\[
\widehat{p}_1
=
\frac{x_1}{n_1}
\]

and

\[
\widehat{p}_2
=
\frac{x_2}{n_2}.
\]

Consider

\[
\boxed{
H_0:p_1=p_2.
}
\]

Under this null, both groups share a common proportion \(p\).

%%%%%%%%%%%%%%%%%%%%%%%%%%%%%%%%%%%%%%%%%%%%%%%%%%%%%%%%%%%%
\subsection{Pooled Proportion Under the Null}
\label{subsec:pooled-proportion-test}
%%%%%%%%%%%%%%%%%%%%%%%%%%%%%%%%%%%%%%%%%%%%%%%%%%%%%%%%%%%%

Under

\[
H_0:p_1=p_2,
\]

estimate the common null proportion by

\[
\boxed{
\widehat{p}_{pool}
=
\frac{
x_1+x_2
}{
n_1+n_2
}.
}
\]

Then the test statistic is

\[
\boxed{
Z
=
\frac{
\widehat{p}_1-\widehat{p}_2
}{
\sqrt{
\widehat{p}_{pool}
(1-\widehat{p}_{pool})
\left(
\frac1{n_1}
+
\frac1{n_2}
\right)
}
}.
}
\]

%%%%%%%%%%%%%%%%%%%%%%%%%%%%%%%%%%%%%%%%%%%%%%%%%%%%%%%%%%%%
\subsection{Worked Example: Two-Proportion Test}
\label{subsec:worked-two-proportion-test}
%%%%%%%%%%%%%%%%%%%%%%%%%%%%%%%%%%%%%%%%%%%%%%%%%%%%%%%%%%%%

\begin{workedexamplebox}{Comparing Two Success Rates}

Suppose

\[
x_1=120,
\qquad
n_1=200,
\]

and

\[
x_2=120,
\qquad
n_2=250.
\]

Then

\[
\widehat{p}_1=0.60,
\qquad
\widehat{p}_2=0.48.
\]

The pooled proportion is

\[
\widehat{p}_{pool}
=
\frac{
240
}{
450
}.
\]

The test statistic is

\[
z
=
\frac{
0.60-0.48
}{
\sqrt{
\widehat{p}_{pool}
(1-\widehat{p}_{pool})
\left(
1/200+1/250
\right)
}
}.
\]

\begin{answerbox}
\[
\boxed{
z
=
\frac{
0.12
}{
\sqrt{
(240/450)(210/450)
(1/200+1/250)
}
}.
}
\]
\end{answerbox}

\end{workedexamplebox}

%%%%%%%%%%%%%%%%%%%%%%%%%%%%%%%%%%%%%%%%%%%%%%%%%%%%%%%%%%%%
\subsection{Testing a Normal Population Variance}
\label{subsec:test-normal-variance}
%%%%%%%%%%%%%%%%%%%%%%%%%%%%%%%%%%%%%%%%%%%%%%%%%%%%%%%%%%%%

Suppose

\[
X_i
\overset{\text{i.i.d.}}{\sim}
N(\mu,\sigma^2).
\]

To test

\[
\boxed{
H_0:\sigma^2=\sigma_0^2,
}
\]

use

\[
\boxed{
\chi^2
=
\frac{
(n-1)S^2
}{
\sigma_0^2
}.
}
\]

Under \(H_0\),

\[
\boxed{
\chi^2
\sim
\chi^2_{n-1}.
}
\]

%%%%%%%%%%%%%%%%%%%%%%%%%%%%%%%%%%%%%%%%%%%%%%%%%%%%%%%%%%%%
\subsection{Two-Sided Variance Test}
\label{subsec:two-sided-variance-test}
%%%%%%%%%%%%%%%%%%%%%%%%%%%%%%%%%%%%%%%%%%%%%%%%%%%%%%%%%%%%

For

\[
H_A:\sigma^2\neq\sigma_0^2,
\]

evidence against \(H_0\) may occur in either tail of the chi-square
distribution.

Thus reject for sufficiently small or sufficiently large values of

\[
\frac{
(n-1)S^2
}{
\sigma_0^2
}.
\]

%%%%%%%%%%%%%%%%%%%%%%%%%%%%%%%%%%%%%%%%%%%%%%%%%%%%%%%%%%%%
\subsection{Testing a Ratio of Two Normal Variances}
\label{subsec:test-ratio-variances}
%%%%%%%%%%%%%%%%%%%%%%%%%%%%%%%%%%%%%%%%%%%%%%%%%%%%%%%%%%%%

Suppose two independent normal populations have variances

\[
\sigma_1^2
\]

and

\[
\sigma_2^2.
\]

To test

\[
H_0:
\frac{
\sigma_1^2
}{
\sigma_2^2
}
=
c,
\]

use

\[
\boxed{
F
=
\frac{
S_1^2/S_2^2
}{
c
}.
}
\]

Under the null,

\[
\boxed{
F
\sim
F_{n_1-1,n_2-1}.
}
\]

For the common hypothesis

\[
H_0:\sigma_1^2=\sigma_2^2,
\]

we have

\[
\boxed{
F
=
\frac{
S_1^2
}{
S_2^2
}.
}
\]

%%%%%%%%%%%%%%%%%%%%%%%%%%%%%%%%%%%%%%%%%%%%%%%%%%%%%%%%%%%%
\subsection{Sensitivity of Variance Tests}
\label{subsec:variance-tests-normality}
%%%%%%%%%%%%%%%%%%%%%%%%%%%%%%%%%%%%%%%%%%%%%%%%%%%%%%%%%%%%

Classical chi-square and \(F\)-tests for variances are sensitive to
departures from normality.

Heavy tails or skewness can substantially affect their Type I error.

Thus these exact tests should not be treated as distribution-free
procedures.

%%%%%%%%%%%%%%%%%%%%%%%%%%%%%%%%%%%%%%%%%%%%%%%%%%%%%%%%%%%%
\subsection{Goodness-of-Fit Tests}
\label{subsec:goodness-of-fit-tests}
%%%%%%%%%%%%%%%%%%%%%%%%%%%%%%%%%%%%%%%%%%%%%%%%%%%%%%%%%%%%

A goodness-of-fit test evaluates whether observed categorical frequencies
are consistent with a proposed probability model.

Suppose categories are

\[
1,\ldots,k.
\]

Let

\[
O_i
\]

be the observed count in category \(i\), and let

\[
E_i
\]

be the expected count under \(H_0\).

%%%%%%%%%%%%%%%%%%%%%%%%%%%%%%%%%%%%%%%%%%%%%%%%%%%%%%%%%%%%
\subsection{Pearson Chi-Square Statistic}
\label{subsec:pearson-chi-square-statistic}
%%%%%%%%%%%%%%%%%%%%%%%%%%%%%%%%%%%%%%%%%%%%%%%%%%%%%%%%%%%%

The Pearson chi-square statistic is

\[
\boxed{
X^2
=
\sum_{i=1}^{k}
\frac{
(O_i-E_i)^2
}{
E_i
}.
}
\]

Large values indicate that observed frequencies differ substantially
from expected frequencies.

Under suitable conditions,

\[
X^2
\]

has an approximate chi-square distribution under the null.

%%%%%%%%%%%%%%%%%%%%%%%%%%%%%%%%%%%%%%%%%%%%%%%%%%%%%%%%%%%%
\subsection{Goodness-of-Fit Expected Counts}
\label{subsec:gof-expected-counts}
%%%%%%%%%%%%%%%%%%%%%%%%%%%%%%%%%%%%%%%%%%%%%%%%%%%%%%%%%%%%

If the null model assigns probability

\[
p_i
\]

to category \(i\), then

\[
\boxed{
E_i
=
np_i.
}
\]

These expected counts, not the observed counts, determine the quality of
the chi-square approximation.

Small expected counts may require combining categories or using exact or
simulation-based methods.

%%%%%%%%%%%%%%%%%%%%%%%%%%%%%%%%%%%%%%%%%%%%%%%%%%%%%%%%%%%%
\subsection{Goodness-of-Fit Degrees of Freedom}
\label{subsec:gof-df}
%%%%%%%%%%%%%%%%%%%%%%%%%%%%%%%%%%%%%%%%%%%%%%%%%%%%%%%%%%%%

If all null probabilities are fully specified, the usual degrees of
freedom are

\[
\boxed{
k-1.
}
\]

If \(m\) parameters are estimated from the same data to obtain the
expected probabilities, a common adjusted formula is

\[
\boxed{
k-1-m,
}
\]

under standard regularity conditions.

%%%%%%%%%%%%%%%%%%%%%%%%%%%%%%%%%%%%%%%%%%%%%%%%%%%%%%%%%%%%
\subsection{Worked Example: Goodness-of-Fit}
\label{subsec:worked-goodness-of-fit}
%%%%%%%%%%%%%%%%%%%%%%%%%%%%%%%%%%%%%%%%%%%%%%%%%%%%%%%%%%%%

\begin{workedexamplebox}{Testing a Fair Die}

A die is rolled

\[
n=120
\]

times.

If the die is fair, each face has probability

\[
\frac16.
\]

Therefore every expected count is

\[
E_i
=
120
\left(
\frac16
\right)
=
20.
\]

If the observed counts are

\[
18,\ 21,\ 16,\ 25,\ 22,\ 18,
\]

then

\[
X^2
=
\frac{(18-20)^2}{20}
+
\frac{(21-20)^2}{20}
+
\cdots
+
\frac{(18-20)^2}{20}.
\]

\begin{answerbox}
\[
\boxed{
X^2
=
\sum_{i=1}^{6}
\frac{
(O_i-20)^2
}{
20
},
\qquad
df=5.
}
\]
\end{answerbox}

\end{workedexamplebox}

%%%%%%%%%%%%%%%%%%%%%%%%%%%%%%%%%%%%%%%%%%%%%%%%%%%%%%%%%%%%
\subsection{Chi-Square Test of Independence}
\label{subsec:chi-square-independence}
%%%%%%%%%%%%%%%%%%%%%%%%%%%%%%%%%%%%%%%%%%%%%%%%%%%%%%%%%%%%

Suppose two categorical variables are recorded in an

\[
r\times c
\]

contingency table.

The null hypothesis is

\[
\boxed{
H_0:
\text{the two variables are independent}.
}
\]

The alternative is that they are associated.

%%%%%%%%%%%%%%%%%%%%%%%%%%%%%%%%%%%%%%%%%%%%%%%%%%%%%%%%%%%%
\subsection{Expected Counts Under Independence}
\label{subsec:expected-counts-independence}
%%%%%%%%%%%%%%%%%%%%%%%%%%%%%%%%%%%%%%%%%%%%%%%%%%%%%%%%%%%%

Let

\[
n_{i\cdot}
\]

be row \(i\)'s total,

\[
n_{\cdot j}
\]

be column \(j\)'s total, and

\[
n
\]

be the grand total.

Under independence,

\[
\boxed{
E_{ij}
=
\frac{
n_{i\cdot}
n_{\cdot j}
}{
n
}.
}
\]

This follows from

\[
P(A_i\cap B_j)
=
P(A_i)P(B_j).
\]

%%%%%%%%%%%%%%%%%%%%%%%%%%%%%%%%%%%%%%%%%%%%%%%%%%%%%%%%%%%%
\subsection{Chi-Square Independence Statistic}
\label{subsec:chi-square-independence-statistic}
%%%%%%%%%%%%%%%%%%%%%%%%%%%%%%%%%%%%%%%%%%%%%%%%%%%%%%%%%%%%

The statistic is

\[
\boxed{
X^2
=
\sum_{i=1}^{r}
\sum_{j=1}^{c}
\frac{
(O_{ij}-E_{ij})^2
}{
E_{ij}
}.
}
\]

Under the null and suitable large-sample conditions,

\[
\boxed{
X^2
\approx
\chi^2_{(r-1)(c-1)}.
}
\]

%%%%%%%%%%%%%%%%%%%%%%%%%%%%%%%%%%%%%%%%%%%%%%%%%%%%%%%%%%%%
\subsection{Worked Example: Expected Cell Count}
\label{subsec:worked-expected-cell-count}
%%%%%%%%%%%%%%%%%%%%%%%%%%%%%%%%%%%%%%%%%%%%%%%%%%%%%%%%%%%%

\begin{workedexamplebox}{A Contingency Table Cell}

Suppose a table has grand total

\[
n=200.
\]

One row total is

\[
80
\]

and one column total is

\[
50.
\]

Under independence, the expected count in their intersection is

\[
E
=
\frac{
80(50)
}{
200
}.
\]

Therefore,

\begin{answerbox}
\[
\boxed{
E=20.
}
\]
\end{answerbox}

\end{workedexamplebox}

%%%%%%%%%%%%%%%%%%%%%%%%%%%%%%%%%%%%%%%%%%%%%%%%%%%%%%%%%%%%
\subsection{Test of Homogeneity}
\label{subsec:test-of-homogeneity}
%%%%%%%%%%%%%%%%%%%%%%%%%%%%%%%%%%%%%%%%%%%%%%%%%%%%%%%%%%%%

A chi-square test of homogeneity compares categorical distributions
across several populations or groups.

Computationally, its Pearson statistic has the same form as the test of
independence.

The distinction is primarily in the study design and interpretation.

Independence asks whether two variables are associated in one sampled
population.

Homogeneity asks whether several groups share the same categorical
distribution.

%%%%%%%%%%%%%%%%%%%%%%%%%%%%%%%%%%%%%%%%%%%%%%%%%%%%%%%%%%%%
\subsection{Residuals in Contingency Tables}
\label{subsec:contingency-residuals}
%%%%%%%%%%%%%%%%%%%%%%%%%%%%%%%%%%%%%%%%%%%%%%%%%%%%%%%%%%%%

A significant chi-square statistic indicates some departure from the
null, but it does not identify which cells are responsible.

A simple Pearson residual is

\[
\boxed{
r_{ij}
=
\frac{
O_{ij}-E_{ij}
}{
\sqrt{E_{ij}}
}.
}
\]

Large residual magnitudes indicate cells contributing strongly to the
overall chi-square statistic.

%%%%%%%%%%%%%%%%%%%%%%%%%%%%%%%%%%%%%%%%%%%%%%%%%%%%%%%%%%%%
\subsection{Likelihood-Ratio Test}
\label{subsec:likelihood-ratio-test}
%%%%%%%%%%%%%%%%%%%%%%%%%%%%%%%%%%%%%%%%%%%%%%%%%%%%%%%%%%%%

Suppose a full model allows

\[
\theta\in\Theta,
\]

while the null restricts

\[
\theta\in\Theta_0.
\]

Let

\[
\widehat{\theta}
\]

be the unrestricted MLE and

\[
\widehat{\theta}_0
\]

the MLE under the null.

Define the likelihood ratio

\[
\boxed{
\Lambda
=
\frac{
L(\widehat{\theta}_0)
}{
L(\widehat{\theta})
}.
}
\]

Since the denominator is the unrestricted maximum,

\[
\boxed{
0\leq\Lambda\leq1.
}
\]

%%%%%%%%%%%%%%%%%%%%%%%%%%%%%%%%%%%%%%%%%%%%%%%%%%%%%%%%%%%%
\subsection{Likelihood-Ratio Statistic}
\label{subsec:likelihood-ratio-statistic}
%%%%%%%%%%%%%%%%%%%%%%%%%%%%%%%%%%%%%%%%%%%%%%%%%%%%%%%%%%%%

A standard transformed statistic is

\[
\boxed{
-2\log\Lambda
=
2
\left[
\ell(\widehat{\theta})
-
\ell(\widehat{\theta}_0)
\right].
}
\]

Large values indicate that the unrestricted model fits much better than
the null-restricted model.

%%%%%%%%%%%%%%%%%%%%%%%%%%%%%%%%%%%%%%%%%%%%%%%%%%%%%%%%%%%%
\subsection{Wilks' Theorem Preview}
\label{subsec:wilks-theorem}
%%%%%%%%%%%%%%%%%%%%%%%%%%%%%%%%%%%%%%%%%%%%%%%%%%%%%%%%%%%%

Under standard regularity conditions, if the null imposes \(q\)
independent restrictions,

\[
\boxed{
-2\log\Lambda
\xrightarrow{d}
\chi^2_q
}
\]

under \(H_0\).

This is Wilks' theorem.

It provides a general asymptotic testing method for parametric models.

%%%%%%%%%%%%%%%%%%%%%%%%%%%%%%%%%%%%%%%%%%%%%%%%%%%%%%%%%%%%
\subsection{Wald Test}
\label{subsec:wald-test}
%%%%%%%%%%%%%%%%%%%%%%%%%%%%%%%%%%%%%%%%%%%%%%%%%%%%%%%%%%%%

Suppose

\[
\widehat{\theta}
\]

is approximately normal:

\[
\widehat{\theta}
\approx
N
(
\theta,
\operatorname{SE}^2
).
\]

To test

\[
H_0:\theta=\theta_0,
\]

a Wald statistic is

\[
\boxed{
Z_W
=
\frac{
\widehat{\theta}-\theta_0
}{
\widehat{\operatorname{SE}}(\widehat{\theta})
}.
}
\]

Under regularity conditions, this is approximately standard normal.

%%%%%%%%%%%%%%%%%%%%%%%%%%%%%%%%%%%%%%%%%%%%%%%%%%%%%%%%%%%%
\subsection{Score Test}
\label{subsec:score-test}
%%%%%%%%%%%%%%%%%%%%%%%%%%%%%%%%%%%%%%%%%%%%%%%%%%%%%%%%%%%%

The score test evaluates the slope of the log-likelihood at the null
parameter value.

For a scalar parameter,

\[
\boxed{
U(\theta_0)
=
\ell'(\theta_0).
}
\]

A standardized score statistic is based on

\[
\boxed{
\frac{
U(\theta_0)
}{
\sqrt{
I(\theta_0)
}
}.
}
\]

The score test does not require fitting the unrestricted model in its
simplest form.

%%%%%%%%%%%%%%%%%%%%%%%%%%%%%%%%%%%%%%%%%%%%%%%%%%%%%%%%%%%%
\subsection{Wald, Score, and Likelihood-Ratio Tests}
\label{subsec:wald-score-lr}
%%%%%%%%%%%%%%%%%%%%%%%%%%%%%%%%%%%%%%%%%%%%%%%%%%%%%%%%%%%%

In regular large-sample problems, the three classical tests are
asymptotically equivalent:

\[
\boxed{
\text{Wald},
}
\]

\[
\boxed{
\text{Score},
}
\]

and

\[
\boxed{
\text{Likelihood Ratio}.
}
\]

However, their finite-sample behavior may differ.

The likelihood-ratio and score procedures often behave better near
parameter boundaries than naive Wald approximations.

%%%%%%%%%%%%%%%%%%%%%%%%%%%%%%%%%%%%%%%%%%%%%%%%%%%%%%%%%%%%
\subsection{Neyman--Pearson Framework}
\label{subsec:neyman-pearson}
%%%%%%%%%%%%%%%%%%%%%%%%%%%%%%%%%%%%%%%%%%%%%%%%%%%%%%%%%%%%

The Neyman--Pearson framework considers testing a simple null hypothesis

\[
H_0:\theta=\theta_0
\]

against a simple alternative

\[
H_1:\theta=\theta_1.
\]

Among all tests with a fixed Type I error probability, one seeks the test
with greatest power against \(\theta_1\).

%%%%%%%%%%%%%%%%%%%%%%%%%%%%%%%%%%%%%%%%%%%%%%%%%%%%%%%%%%%%
\subsection{Neyman--Pearson Lemma}
\label{subsec:neyman-pearson-lemma}
%%%%%%%%%%%%%%%%%%%%%%%%%%%%%%%%%%%%%%%%%%%%%%%%%%%%%%%%%%%%

\begin{theorem}[Neyman--Pearson Lemma]
For testing a simple null hypothesis against a simple alternative, a
most powerful level-\(\alpha\) test rejects \(H_0\) for sufficiently
large values of the likelihood ratio

\[
\boxed{
\frac{
L(\theta_1;x)
}{
L(\theta_0;x)
}.
}
\]
\end{theorem}

Thus likelihood ratios arise naturally from optimal testing theory.

%%%%%%%%%%%%%%%%%%%%%%%%%%%%%%%%%%%%%%%%%%%%%%%%%%%%%%%%%%%%
\subsection{Most Powerful Tests}
\label{subsec:most-powerful-tests}
%%%%%%%%%%%%%%%%%%%%%%%%%%%%%%%%%%%%%%%%%%%%%%%%%%%%%%%%%%%%

A test is most powerful at a specified alternative if it has the largest
power among all tests of the same significance level.

Thus, among level-\(\alpha\) tests,

\[
\boxed{
\text{best test}
=
\text{largest probability of rejecting }H_0
\text{ under the chosen alternative}.
}
\]

%%%%%%%%%%%%%%%%%%%%%%%%%%%%%%%%%%%%%%%%%%%%%%%%%%%%%%%%%%%%
\subsection{Uniformly Most Powerful Tests}
\label{subsec:ump-tests}
%%%%%%%%%%%%%%%%%%%%%%%%%%%%%%%%%%%%%%%%%%%%%%%%%%%%%%%%%%%%

A uniformly most powerful, or UMP, test has maximum power simultaneously
for every parameter value in a composite alternative.

Such tests do not always exist.

They arise in important one-parameter families with monotone likelihood
ratio structure.

%%%%%%%%%%%%%%%%%%%%%%%%%%%%%%%%%%%%%%%%%%%%%%%%%%%%%%%%%%%%
\subsection{Power Function}
\label{subsec:power-function}
%%%%%%%%%%%%%%%%%%%%%%%%%%%%%%%%%%%%%%%%%%%%%%%%%%%%%%%%%%%%

For a rejection region \(R\), define

\[
\boxed{
\pi(\theta)
=
P_\theta(T\in R).
}
\]

This function is called the power function.

For \(\theta\) satisfying the null hypothesis, the power should be at
most the test size.

For alternative parameter values, larger power is preferred.

%%%%%%%%%%%%%%%%%%%%%%%%%%%%%%%%%%%%%%%%%%%%%%%%%%%%%%%%%%%%
\subsection{Power of a One-Sided \(z\)-Test}
\label{subsec:power-one-sided-z}
%%%%%%%%%%%%%%%%%%%%%%%%%%%%%%%%%%%%%%%%%%%%%%%%%%%%%%%%%%%%

Suppose

\[
X_i
\sim
N(\mu,\sigma^2)
\]

with known \(\sigma\).

Test

\[
H_0:\mu=\mu_0
\]

against

\[
H_A:\mu>\mu_0.
\]

Reject when

\[
\overline{X}
>
\mu_0
+
z_\alpha
\frac{\sigma}{\sqrt{n}}.
\]

If the true mean is \(\mu_1>\mu_0\), then power is

\[
\boxed{
P_{\mu_1}
\left(
\overline{X}
>
\mu_0
+
z_\alpha
\frac{\sigma}{\sqrt{n}}
\right).
}
\]

%%%%%%%%%%%%%%%%%%%%%%%%%%%%%%%%%%%%%%%%%%%%%%%%%%%%%%%%%%%%
\subsection{Power Formula for the One-Sided \(z\)-Test}
\label{subsec:power-formula-z}
%%%%%%%%%%%%%%%%%%%%%%%%%%%%%%%%%%%%%%%%%%%%%%%%%%%%%%%%%%%%

Under \(\mu=\mu_1\),

\[
\overline{X}
\sim
N
\left(
\mu_1,
\frac{\sigma^2}{n}
\right).
\]

Therefore,

\[
\boxed{
\operatorname{Power}(\mu_1)
=
1-
\Phi
\left[
z_\alpha
-
\frac{
(\mu_1-\mu_0)\sqrt{n}
}{
\sigma
}
\right].
}
\]

Power increases as the standardized effect size increases.

%%%%%%%%%%%%%%%%%%%%%%%%%%%%%%%%%%%%%%%%%%%%%%%%%%%%%%%%%%%%
\subsection{Factors Affecting Power}
\label{subsec:factors-affecting-power}
%%%%%%%%%%%%%%%%%%%%%%%%%%%%%%%%%%%%%%%%%%%%%%%%%%%%%%%%%%%%

Power generally increases when:

\[
\boxed{
\text{sample size increases},
}
\]

\[
\boxed{
\text{effect size increases},
}
\]

\[
\boxed{
\text{measurement variability decreases},
}
\]

or

\[
\boxed{
\alpha\text{ increases}.
}
\]

However, increasing \(\alpha\) also increases Type I error risk.

%%%%%%%%%%%%%%%%%%%%%%%%%%%%%%%%%%%%%%%%%%%%%%%%%%%%%%%%%%%%
\subsection{Effect Size}
\label{subsec:effect-size-testing}
%%%%%%%%%%%%%%%%%%%%%%%%%%%%%%%%%%%%%%%%%%%%%%%%%%%%%%%%%%%%

For a one-sample normal mean, a standardized effect size is

\[
\boxed{
d
=
\frac{
\mu_1-\mu_0
}{
\sigma
}.
}
\]

Then the power expression depends on

\[
d\sqrt{n}.
\]

Thus a small effect can become detectable with sufficiently large
sample size.

%%%%%%%%%%%%%%%%%%%%%%%%%%%%%%%%%%%%%%%%%%%%%%%%%%%%%%%%%%%%
\subsection{Cohen's \(d\) Preview}
\label{subsec:cohens-d-preview}
%%%%%%%%%%%%%%%%%%%%%%%%%%%%%%%%%%%%%%%%%%%%%%%%%%%%%%%%%%%%

For observed one-sample data, a standardized mean difference may be
estimated by

\[
\boxed{
d
=
\frac{
\overline{x}-\mu_0
}{
s
}.
}
\]

For two independent groups, standardized mean differences use a suitable
pooled or otherwise justified scale estimate.

Effect-size interpretation should depend on subject-matter context rather
than rigid universal cutoffs.

%%%%%%%%%%%%%%%%%%%%%%%%%%%%%%%%%%%%%%%%%%%%%%%%%%%%%%%%%%%%
\subsection{Sample Size for Testing a Mean}
\label{subsec:sample-size-testing-mean}
%%%%%%%%%%%%%%%%%%%%%%%%%%%%%%%%%%%%%%%%%%%%%%%%%%%%%%%%%%%%

For a one-sided \(z\)-test of a mean with known \(\sigma\), desired
significance level \(\alpha\), power \(1-\beta\), and effect

\[
\Delta
=
\mu_1-\mu_0>0,
\]

a common planning formula is

\[
\boxed{
n
=
\left[
\frac{
(z_\alpha+z_\beta)\sigma
}{
\Delta
}
\right]^2,
}
\]

where the quantiles are interpreted as positive upper critical values.

Round upward.

%%%%%%%%%%%%%%%%%%%%%%%%%%%%%%%%%%%%%%%%%%%%%%%%%%%%%%%%%%%%
\subsection{Worked Example: Power-Based Sample Size}
\label{subsec:worked-power-sample-size}
%%%%%%%%%%%%%%%%%%%%%%%%%%%%%%%%%%%%%%%%%%%%%%%%%%%%%%%%%%%%

\begin{workedexamplebox}{Planning a One-Sided Mean Test}

Suppose

\[
\sigma=10,
\]

a difference of

\[
\Delta=5
\]

should be detected,

\[
\alpha=0.05,
\]

and desired power is

\[
0.80.
\]

Using approximately

\[
z_{0.05}=1.645
\]

and

\[
z_{0.20}=0.842,
\]

we obtain

\[
n
=
\left[
\frac{
(1.645+0.842)(10)
}{
5
}
\right]^2.
\]

Therefore,

\begin{answerbox}
\[
\boxed{
n
\approx
24.74,
}
\]

so use at least

\[
\boxed{
25
}
\]

observations under this planning approximation.
\end{answerbox}

\end{workedexamplebox}

%%%%%%%%%%%%%%%%%%%%%%%%%%%%%%%%%%%%%%%%%%%%%%%%%%%%%%%%%%%%
\subsection{Two-Sided Tests Require More Evidence in Each Tail}
\label{subsec:two-sided-power}
%%%%%%%%%%%%%%%%%%%%%%%%%%%%%%%%%%%%%%%%%%%%%%%%%%%%%%%%%%%%

For a two-sided level-\(\alpha\) test, each tail receives approximately

\[
\frac{\alpha}{2}.
\]

The critical value is therefore larger than in the corresponding
one-sided test.

For the same sample size and directional effect, a two-sided test usually
has less power than a correctly pre-specified one-sided test.

%%%%%%%%%%%%%%%%%%%%%%%%%%%%%%%%%%%%%%%%%%%%%%%%%%%%%%%%%%%%
\subsection{When One-Sided Tests Are Appropriate}
\label{subsec:when-one-sided}
%%%%%%%%%%%%%%%%%%%%%%%%%%%%%%%%%%%%%%%%%%%%%%%%%%%%%%%%%%%%

A one-sided test should be chosen because the scientific hypothesis is
directional before inspecting the data.

It should not be chosen after observing which direction produced a more
favorable result.

Changing the direction after seeing the data invalidates the nominal
Type I error control.

%%%%%%%%%%%%%%%%%%%%%%%%%%%%%%%%%%%%%%%%%%%%%%%%%%%%%%%%%%%%
\subsection{Confidence Interval Duality}
\label{subsec:testing-confidence-duality}
%%%%%%%%%%%%%%%%%%%%%%%%%%%%%%%%%%%%%%%%%%%%%%%%%%%%%%%%%%%%

For many standard procedures, a two-sided level-\(\alpha\) test of

\[
H_0:\theta=\theta_0
\]

rejects exactly when

\[
\theta_0
\]

lies outside the corresponding

\[
100(1-\alpha)\%
\]

confidence interval.

Thus

\[
\boxed{
\text{tests}
\leftrightarrow
\text{confidence intervals}.
}
\]

Confidence intervals usually provide more information because they
display the range of parameter values compatible with the data.

%%%%%%%%%%%%%%%%%%%%%%%%%%%%%%%%%%%%%%%%%%%%%%%%%%%%%%%%%%%%
\subsection{Equivalence Testing}
\label{subsec:equivalence-testing-preview}
%%%%%%%%%%%%%%%%%%%%%%%%%%%%%%%%%%%%%%%%%%%%%%%%%%%%%%%%%%%%

Sometimes the scientific question is not whether two quantities differ,
but whether they are sufficiently similar.

Suppose differences within

\[
(-\Delta,\Delta)
\]

are considered practically equivalent.

An equivalence test reverses the usual roles conceptually:

\[
\boxed{
H_0:
|\theta|\geq\Delta
}
\]

against

\[
\boxed{
H_A:
|\theta|<\Delta.
}
\]

This is not the same as failing to reject a zero-difference null.

%%%%%%%%%%%%%%%%%%%%%%%%%%%%%%%%%%%%%%%%%%%%%%%%%%%%%%%%%%%%
\subsection{TOST Procedure Preview}
\label{subsec:tost-preview}
%%%%%%%%%%%%%%%%%%%%%%%%%%%%%%%%%%%%%%%%%%%%%%%%%%%%%%%%%%%%

The Two One-Sided Tests procedure, abbreviated TOST, tests

\[
H_{01}:
\theta\leq-\Delta
\]

and

\[
H_{02}:
\theta\geq\Delta.
\]

Equivalence is concluded only if both null hypotheses are rejected.

This corresponds closely to checking whether an appropriately chosen
confidence interval lies entirely within the equivalence region.

%%%%%%%%%%%%%%%%%%%%%%%%%%%%%%%%%%%%%%%%%%%%%%%%%%%%%%%%%%%%
\subsection{Noninferiority Testing Preview}
\label{subsec:noninferiority-preview}
%%%%%%%%%%%%%%%%%%%%%%%%%%%%%%%%%%%%%%%%%%%%%%%%%%%%%%%%%%%%

A noninferiority test asks whether a new treatment is not worse than a
reference treatment by more than a prespecified margin.

For effect

\[
\theta
=
\mu_{\text{new}}
-
\mu_{\text{reference}},
\]

one possible formulation is

\[
H_0:
\theta\leq-\Delta
\]

against

\[
H_A:
\theta>-\Delta.
\]

The sign convention depends on how favorable outcomes are defined.

%%%%%%%%%%%%%%%%%%%%%%%%%%%%%%%%%%%%%%%%%%%%%%%%%%%%%%%%%%%%
\subsection{Permutation Tests}
\label{subsec:permutation-tests}
%%%%%%%%%%%%%%%%%%%%%%%%%%%%%%%%%%%%%%%%%%%%%%%%%%%%%%%%%%%%

Permutation tests use the distribution generated by rearranging group
labels under a null hypothesis of exchangeability.

Suppose two groups are compared using statistic

\[
T.
\]

Under a randomization null:

\begin{enumerate}
    \item pool observations;
    \item randomly permute treatment labels;
    \item recompute \(T\);
    \item repeat many times or enumerate all valid assignments;
    \item compare the observed statistic with the permutation
          distribution.
\end{enumerate}

%%%%%%%%%%%%%%%%%%%%%%%%%%%%%%%%%%%%%%%%%%%%%%%%%%%%%%%%%%%%
\subsection{Randomization Tests}
\label{subsec:randomization-tests}
%%%%%%%%%%%%%%%%%%%%%%%%%%%%%%%%%%%%%%%%%%%%%%%%%%%%%%%%%%%%

In randomized experiments, the treatment assignment mechanism itself can
generate a null distribution.

This provides a design-based test that does not necessarily require a
parametric population model.

The validity comes from the known randomization procedure.

%%%%%%%%%%%%%%%%%%%%%%%%%%%%%%%%%%%%%%%%%%%%%%%%%%%%%%%%%%%%
\subsection{Worked Example: Permutation Logic}
\label{subsec:worked-permutation-logic}
%%%%%%%%%%%%%%%%%%%%%%%%%%%%%%%%%%%%%%%%%%%%%%%%%%%%%%%%%%%%

\begin{workedexamplebox}{Difference in Means}

Suppose treatment labels were randomly assigned.

The observed statistic is

\[
T_{\text{obs}}
=
\overline{x}_{T}
-
\overline{x}_{C}.
\]

Under a sharp no-effect null, treatment labels may be reassigned according
to the original randomization design.

Each reassignment produces

\[
T^*.
\]

The \(p\)-value is estimated by the proportion of randomizations satisfying

\[
|T^*|
\geq
|T_{\text{obs}}|
\]

for a two-sided test.

\begin{answerbox}
\[
\boxed{
p
\approx
\frac{
\#\{
|T^*|\geq|T_{\text{obs}}|
\}
}{
\text{number of randomizations}
}.
}
\]
\end{answerbox}

\end{workedexamplebox}

%%%%%%%%%%%%%%%%%%%%%%%%%%%%%%%%%%%%%%%%%%%%%%%%%%%%%%%%%%%%
\subsection{Permutation Assumptions}
\label{subsec:permutation-assumptions}
%%%%%%%%%%%%%%%%%%%%%%%%%%%%%%%%%%%%%%%%%%%%%%%%%%%%%%%%%%%%

Permutation tests are not automatically assumption-free.

Their validity depends on an exchangeability or randomization structure.

For independent observational groups with different distributions or
variances, indiscriminate permutation may not preserve the intended null
distribution.

The permutation scheme must reflect the scientific design.

%%%%%%%%%%%%%%%%%%%%%%%%%%%%%%%%%%%%%%%%%%%%%%%%%%%%%%%%%%%%
\subsection{Multiple Testing}
\label{subsec:multiple-testing}
%%%%%%%%%%%%%%%%%%%%%%%%%%%%%%%%%%%%%%%%%%%%%%%%%%%%%%%%%%%%

Suppose many hypotheses are tested simultaneously.

Even if every null is true, testing each at level

\[
\alpha=0.05
\]

creates repeated opportunities for false rejection.

For \(m\) independent tests, the probability of at least one false
rejection is

\[
\boxed{
1-(1-\alpha)^m.
}
\]

This can become large when \(m\) is large.

%%%%%%%%%%%%%%%%%%%%%%%%%%%%%%%%%%%%%%%%%%%%%%%%%%%%%%%%%%%%
\subsection{Family-Wise Error Rate}
\label{subsec:family-wise-error-rate}
%%%%%%%%%%%%%%%%%%%%%%%%%%%%%%%%%%%%%%%%%%%%%%%%%%%%%%%%%%%%

\begin{definitionbox}{Family-Wise Error Rate}
The family-wise error rate, abbreviated FWER, is the probability of
making at least one Type I error among a family of hypothesis tests.
\end{definitionbox}

Symbolically,

\[
\boxed{
FWER
=
P(V\geq1),
}
\]

where \(V\) is the number of false rejections.

%%%%%%%%%%%%%%%%%%%%%%%%%%%%%%%%%%%%%%%%%%%%%%%%%%%%%%%%%%%%
\subsection{Bonferroni Correction}
\label{subsec:bonferroni-testing}
%%%%%%%%%%%%%%%%%%%%%%%%%%%%%%%%%%%%%%%%%%%%%%%%%%%%%%%%%%%%

Suppose \(m\) hypotheses are tested and the desired overall FWER is at
most \(\alpha\).

Bonferroni tests each individual hypothesis at level

\[
\boxed{
\frac{\alpha}{m}.
}
\]

By the union bound,

\[
\boxed{
P(\text{at least one false rejection})
\leq
m
\frac{\alpha}{m}
=
\alpha.
}
\]

No independence assumption is required for this bound.

%%%%%%%%%%%%%%%%%%%%%%%%%%%%%%%%%%%%%%%%%%%%%%%%%%%%%%%%%%%%
\subsection{Worked Example: Bonferroni}
\label{subsec:worked-bonferroni}
%%%%%%%%%%%%%%%%%%%%%%%%%%%%%%%%%%%%%%%%%%%%%%%%%%%%%%%%%%%%

\begin{workedexamplebox}{Ten Simultaneous Tests}

Suppose

\[
m=10
\]

tests are performed with desired family-wise error rate

\[
0.05.
\]

Bonferroni uses individual significance level

\[
\frac{0.05}{10}
=
0.005.
\]

\begin{answerbox}
\[
\boxed{
\alpha_{\text{individual}}
=
0.005.
}
\]
\end{answerbox}

\end{workedexamplebox}

%%%%%%%%%%%%%%%%%%%%%%%%%%%%%%%%%%%%%%%%%%%%%%%%%%%%%%%%%%%%
\subsection{Holm Procedure Preview}
\label{subsec:holm-procedure-preview}
%%%%%%%%%%%%%%%%%%%%%%%%%%%%%%%%%%%%%%%%%%%%%%%%%%%%%%%%%%%%

Holm's step-down procedure improves on simple Bonferroni while still
controlling FWER.

Order the \(p\)-values:

\[
p_{(1)}
\leq
p_{(2)}
\leq
\cdots
\leq
p_{(m)}.
\]

Compare sequentially against thresholds

\[
\boxed{
\frac{\alpha}{m},
\frac{\alpha}{m-1},
\ldots,
\alpha.
}
\]

Holm's method is uniformly at least as powerful as simple Bonferroni.

%%%%%%%%%%%%%%%%%%%%%%%%%%%%%%%%%%%%%%%%%%%%%%%%%%%%%%%%%%%%
\subsection{False Discovery Rate}
\label{subsec:false-discovery-rate}
%%%%%%%%%%%%%%%%%%%%%%%%%%%%%%%%%%%%%%%%%%%%%%%%%%%%%%%%%%%%

When thousands of hypotheses are tested, strict FWER control may be too
conservative.

The false discovery rate, abbreviated FDR, focuses on the expected
proportion of false discoveries among rejected hypotheses.

If

\[
R
\]

is the total number of rejections and

\[
V
\]

the number of false rejections, then conceptually

\[
\boxed{
FDR
=
\mathbb{E}
\left[
\frac{
V
}{
\max(R,1)
}
\right].
}
\]

%%%%%%%%%%%%%%%%%%%%%%%%%%%%%%%%%%%%%%%%%%%%%%%%%%%%%%%%%%%%
\subsection{Benjamini--Hochberg Procedure Preview}
\label{subsec:benjamini-hochberg-preview}
%%%%%%%%%%%%%%%%%%%%%%%%%%%%%%%%%%%%%%%%%%%%%%%%%%%%%%%%%%%%

Order \(m\) \(p\)-values:

\[
p_{(1)}
\leq
\cdots
\leq
p_{(m)}.
\]

For desired FDR level \(q\), find the largest \(k\) satisfying

\[
\boxed{
p_{(k)}
\leq
\frac{k}{m}q.
}
\]

Then reject the hypotheses associated with

\[
p_{(1)},
\ldots,
p_{(k)}.
\]

Under suitable dependence conditions, this controls FDR.

%%%%%%%%%%%%%%%%%%%%%%%%%%%%%%%%%%%%%%%%%%%%%%%%%%%%%%%%%%%%
\subsection{Data Snooping and Selective Inference}
\label{subsec:data-snooping}
%%%%%%%%%%%%%%%%%%%%%%%%%%%%%%%%%%%%%%%%%%%%%%%%%%%%%%%%%%%%

If many analyses are attempted but only the most favorable result is
reported, nominal \(p\)-values no longer reflect the true analysis
process.

Examples include:

\[
\boxed{
\text{testing many outcomes},
}
\]

\[
\boxed{
\text{trying many subgroups},
}
\]

\[
\boxed{
\text{trying many models},
}
\]

and

\[
\boxed{
\text{stopping data collection when significance appears}.
}
\]

The inferential procedure must account for the full selection process.

%%%%%%%%%%%%%%%%%%%%%%%%%%%%%%%%%%%%%%%%%%%%%%%%%%%%%%%%%%%%
\subsection{Optional Stopping}
\label{subsec:optional-stopping-testing}
%%%%%%%%%%%%%%%%%%%%%%%%%%%%%%%%%%%%%%%%%%%%%%%%%%%%%%%%%%%%

Repeatedly checking a conventional \(p\)-value during data collection and
stopping when

\[
p<0.05
\]

can inflate Type I error.

Sequential testing requires methods designed for repeated looks at the
data.

Examples include:

\[
\boxed{
\text{group-sequential designs},
}
\]

and

\[
\boxed{
\text{alpha-spending procedures}.
}
\]

%%%%%%%%%%%%%%%%%%%%%%%%%%%%%%%%%%%%%%%%%%%%%%%%%%%%%%%%%%%%
\subsection{Sequential Probability Ratio Test Preview}
\label{subsec:sprt-preview}
%%%%%%%%%%%%%%%%%%%%%%%%%%%%%%%%%%%%%%%%%%%%%%%%%%%%%%%%%%%%

The sequential probability ratio test, developed from
Neyman--Pearson ideas, evaluates evidence as observations accumulate.

For simple hypotheses,

\[
H_0:\theta=\theta_0,
\qquad
H_1:\theta=\theta_1,
\]

one monitors the likelihood ratio

\[
\boxed{
\frac{
L_n(\theta_1)
}{
L_n(\theta_0)
}.
}
\]

Sampling may stop when this ratio crosses a predetermined boundary.

%%%%%%%%%%%%%%%%%%%%%%%%%%%%%%%%%%%%%%%%%%%%%%%%%%%%%%%%%%%%
\subsection{Testing Workflow}
\label{subsec:hypothesis-testing-workflow}
%%%%%%%%%%%%%%%%%%%%%%%%%%%%%%%%%%%%%%%%%%%%%%%%%%%%%%%%%%%%

A systematic hypothesis-testing workflow is:

\begin{enumerate}
    \item identify the population parameter;
    \item state \(H_0\);
    \item state \(H_A\);
    \item choose one-sided or two-sided direction before examining the
          result;
    \item identify the sampling design;
    \item select the test statistic;
    \item verify assumptions;
    \item choose significance level \(\alpha\);
    \item compute the observed statistic;
    \item calculate the \(p\)-value or compare with a critical value;
    \item make the reject/fail-to-reject decision;
    \item report an effect estimate and confidence interval;
    \item interpret practical significance and study limitations.
\end{enumerate}

%%%%%%%%%%%%%%%%%%%%%%%%%%%%%%%%%%%%%%%%%%%%%%%%%%%%%%%%%%%%
\subsection{Choosing the Correct Test}
\label{subsec:choosing-correct-test}
%%%%%%%%%%%%%%%%%%%%%%%%%%%%%%%%%%%%%%%%%%%%%%%%%%%%%%%%%%%%

For one quantitative mean with known \(\sigma\):

\[
\boxed{
\text{one-sample }z\text{-test}.
}
\]

For one quantitative mean with unknown \(\sigma\):

\[
\boxed{
\text{one-sample }t\text{-test}.
}
\]

For two independent means:

\[
\boxed{
\text{Welch }t\text{-test}
}
\]

or, under a justified equal-variance model,

\[
\boxed{
\text{pooled }t\text{-test}.
}
\]

For paired observations:

\[
\boxed{
\text{paired }t\text{-test}.
}
\]

For one proportion:

\[
\boxed{
\text{one-proportion }z\text{-test or exact binomial test}.
}
\]

For two proportions:

\[
\boxed{
\text{two-proportion }z\text{-test}.
}
\]

For categorical count tables:

\[
\boxed{
\chi^2\text{ procedures}.
}
\]

%%%%%%%%%%%%%%%%%%%%%%%%%%%%%%%%%%%%%%%%%%%%%%%%%%%%%%%%%%%%
\subsection{Common Errors}
\label{subsec:hypothesis-testing-common-errors}
%%%%%%%%%%%%%%%%%%%%%%%%%%%%%%%%%%%%%%%%%%%%%%%%%%%%%%%%%%%%

\subsubsection{Interpreting the \(p\)-Value as the Probability the Null Is True}

A frequentist \(p\)-value is computed assuming the null is true.

It is not

\[
P(H_0\mid\text{data}).
\]

%%%%%%%%%%%%%%%%%%%%%%%%%%%%%%%%%%%%%%%%%%%%%%%%%%%%%%%%%%%%
\subsubsection{Interpreting \(1-p\) as the Probability the Alternative Is True}

In general,

\[
\boxed{
1-p
}
\]

is not a posterior probability for the alternative hypothesis.

%%%%%%%%%%%%%%%%%%%%%%%%%%%%%%%%%%%%%%%%%%%%%%%%%%%%%%%%%%%%
\subsubsection{Accepting the Null After a Nonsignificant Result}

A nonsignificant result may reflect low power.

Use

\[
\boxed{
\text{fail to reject }H_0
}
\]

unless an equivalence or noninferiority procedure actually supports a
different scientific conclusion.

%%%%%%%%%%%%%%%%%%%%%%%%%%%%%%%%%%%%%%%%%%%%%%%%%%%%%%%%%%%%
\subsubsection{Choosing One-Sided Versus Two-Sided After Seeing the Data}

The direction should be based on the scientific question before
observing the result.

Post hoc direction selection changes the Type I error rate.

%%%%%%%%%%%%%%%%%%%%%%%%%%%%%%%%%%%%%%%%%%%%%%%%%%%%%%%%%%%%
\subsubsection{Using the Estimated Proportion in the One-Proportion Null Standard Error}

For testing

\[
H_0:p=p_0,
\]

the large-sample null standard error is

\[
\boxed{
\sqrt{
\frac{
p_0(1-p_0)
}{n}
}.
}
\]

The confidence-interval standard error often instead uses
\(\widehat{p}\).

%%%%%%%%%%%%%%%%%%%%%%%%%%%%%%%%%%%%%%%%%%%%%%%%%%%%%%%%%%%%
\subsubsection{Forgetting the Pooled Proportion in a Two-Proportion Test}

Under

\[
H_0:p_1=p_2,
\]

the null assumes one common proportion.

Therefore the standard error uses

\[
\boxed{
\widehat{p}_{pool}.
}
\]

%%%%%%%%%%%%%%%%%%%%%%%%%%%%%%%%%%%%%%%%%%%%%%%%%%%%%%%%%%%%
\subsubsection{Pooling Variances Without Justification}

The pooled two-sample \(t\)-test requires a common-variance model.

Welch's test avoids this assumption.

%%%%%%%%%%%%%%%%%%%%%%%%%%%%%%%%%%%%%%%%%%%%%%%%%%%%%%%%%%%%
\subsubsection{Ignoring Pairing}

Paired data must be analyzed through within-pair differences or another
model that explicitly incorporates the dependence.

%%%%%%%%%%%%%%%%%%%%%%%%%%%%%%%%%%%%%%%%%%%%%%%%%%%%%%%%%%%%
\subsubsection{Treating Statistical Significance as Effect Size}

A small \(p\)-value does not measure how large an effect is.

Report the estimated effect and confidence interval.

%%%%%%%%%%%%%%%%%%%%%%%%%%%%%%%%%%%%%%%%%%%%%%%%%%%%%%%%%%%%
\subsubsection{Treating a Large \(p\)-Value as Evidence of Equivalence}

Failure to detect a difference does not establish that the difference is
small enough to be practically negligible.

Use equivalence testing when equivalence is the scientific goal.

%%%%%%%%%%%%%%%%%%%%%%%%%%%%%%%%%%%%%%%%%%%%%%%%%%%%%%%%%%%%
\subsubsection{Ignoring Power}

A test with very low power can fail to detect scientifically important
effects.

Sample-size planning should consider desired power.

%%%%%%%%%%%%%%%%%%%%%%%%%%%%%%%%%%%%%%%%%%%%%%%%%%%%%%%%%%%%
\subsubsection{Ignoring Multiple Testing}

Conducting many unadjusted tests raises the probability of false
discoveries.

The testing family and error criterion should be specified.

%%%%%%%%%%%%%%%%%%%%%%%%%%%%%%%%%%%%%%%%%%%%%%%%%%%%%%%%%%%%
\subsubsection{Using Chi-Square Approximations with Very Small Expected Counts}

The asymptotic chi-square approximation may be poor when expected cell
counts are too small.

Exact, simulation, or alternative methods may be needed.

%%%%%%%%%%%%%%%%%%%%%%%%%%%%%%%%%%%%%%%%%%%%%%%%%%%%%%%%%%%%
\subsubsection{Ignoring Study Design}

A correct test statistic does not repair:

\[
\boxed{
\text{selection bias},
}
\]

\[
\boxed{
\text{confounding},
}
\]

or

\[
\boxed{
\text{invalid measurement}.
}
\]

Statistical significance is only as meaningful as the design that
generated the data.

%%%%%%%%%%%%%%%%%%%%%%%%%%%%%%%%%%%%%%%%%%%%%%%%%%%%%%%%%%%%
\subsection{Applications}
\label{subsec:hypothesis-testing-applications}
%%%%%%%%%%%%%%%%%%%%%%%%%%%%%%%%%%%%%%%%%%%%%%%%%%%%%%%%%%%%

\begin{applicationbox}

\textbf{Scientific Experiments.}
Hypothesis tests compare treatment groups and evaluate whether observed
effects are compatible with null models.

\medskip

\textbf{Medicine.}
Tests evaluate treatment effects, response proportions, risk
differences, and noninferiority claims.

\medskip

\textbf{Engineering.}
Testing procedures evaluate process means, defect rates, variability,
and quality-control claims.

\medskip

\textbf{Environmental Science.}
Researchers test changes in environmental conditions, differences among
sites, and associations between categorical variables.

\medskip

\textbf{Education Research.}
Tests compare instructional methods, group outcomes, survey proportions,
and program effects.

\medskip

\textbf{Business and Economics.}
Statistical tests evaluate changes in demand, conversion rates, economic
parameters, and experimental interventions.

\medskip

\textbf{Data Science.}
Model coefficients, feature effects, A/B experiments, and multiple
comparisons use hypothesis-testing ideas.

\medskip

\textbf{Genomics and High-Dimensional Science.}
Thousands of simultaneous tests require false-discovery-rate methods and
careful multiple-testing control.

\end{applicationbox}

%%%%%%%%%%%%%%%%%%%%%%%%%%%%%%%%%%%%%%%%%%%%%%%%%%%%%%%%%%%%
\subsection{Exercises}
\label{subsec:hypothesis-testing-exercises}
%%%%%%%%%%%%%%%%%%%%%%%%%%%%%%%%%%%%%%%%%%%%%%%%%%%%%%%%%%%%

\begin{exercise}[1]
Define the null hypothesis.
\end{exercise}

\begin{exercise}[1]
Define the alternative hypothesis.
\end{exercise}

\begin{exercise}[1]
Distinguish simple and composite hypotheses.
\end{exercise}

\begin{exercise}[1]
Distinguish one-sided and two-sided tests.
\end{exercise}

\begin{exercise}[1]
Define a Type I error.
\end{exercise}

\begin{exercise}[1]
Define a Type II error.
\end{exercise}

\begin{exercise}[1]
Define test power.
\end{exercise}

\begin{exercise}[1]
Define a \(p\)-value.
\end{exercise}

\begin{exercise}[1]
Explain why a \(p\)-value is not the probability that \(H_0\) is true.
\end{exercise}

\begin{exercise}[1]
Define statistical significance.
\end{exercise}

\begin{exercise}[2]
Suppose

\[
\overline{x}=103,
\qquad
\mu_0=100,
\qquad
\sigma=12,
\qquad
n=64.
\]

Compute the one-sample \(z\)-test statistic.
\end{exercise}

\begin{exercise}[2]
For the previous problem, write the two-sided \(p\)-value using
\(\Phi\).
\end{exercise}

\begin{exercise}[2]
Suppose

\[
n=36,
\qquad
\overline{x}=54,
\qquad
s=12.
\]

Test

\[
H_0:\mu=50.
\]

Compute the \(t\)-statistic and degrees of freedom.
\end{exercise}

\begin{exercise}[2]
Suppose

\[
n=100,
\qquad
x=62.
\]

Test

\[
H_0:p=0.50.
\]

Compute the one-proportion \(z\)-statistic.
\end{exercise}

\begin{exercise}[2]
For two samples,

\[
\overline{x}_1=20,\quad s_1=5,\quad n_1=30,
\]

and

\[
\overline{x}_2=17,\quad s_2=6,\quad n_2=40,
\]

compute the Welch test statistic for

\[
H_0:\mu_1=\mu_2.
\]
\end{exercise}

\begin{exercise}[2]
Suppose paired differences satisfy

\[
\overline{d}=2,
\qquad
s_D=3,
\qquad
n=25.
\]

Compute the paired \(t\)-statistic for

\[
H_0:\mu_D=0.
\]
\end{exercise}

\begin{exercise}[2]
Suppose

\[
x_1=80,\ n_1=100,
\qquad
x_2=60,\ n_2=100.
\]

Compute the pooled proportion under

\[
H_0:p_1=p_2.
\]
\end{exercise}

\begin{exercise}[2]
Use the pooled proportion from the previous problem to write the
two-proportion \(z\)-statistic.
\end{exercise}

\begin{exercise}[2]
Suppose a normal sample has

\[
n=20,
\qquad
s^2=16.
\]

Write the chi-square statistic for testing

\[
H_0:\sigma^2=9.
\]
\end{exercise}

\begin{exercise}[2]
A contingency table cell has row total \(70\), column total \(60\), and
grand total \(210\). Find its expected count under independence.
\end{exercise}

\begin{exercise}[2]
Suppose six categories are completely specified under a
goodness-of-fit null. Find the usual degrees of freedom.
\end{exercise}

\begin{exercise}[2]
If \(20\) tests are Bonferroni-adjusted to control FWER at \(0.05\),
find the individual significance level.
\end{exercise}

\begin{exercise}[3]
Derive the one-sample \(z\)-test rejection region for a two-sided
alternative.
\end{exercise}

\begin{exercise}[3]
Derive the one-proportion null standard error.
\end{exercise}

\begin{exercise}[3]
Explain why the one-proportion hypothesis test uses \(p_0\) in the
standard error while a Wald confidence interval uses \(\widehat{p}\).
\end{exercise}

\begin{exercise}[3]
Derive the Welch standard error for two independent means.
\end{exercise}

\begin{exercise}[3]
Derive the pooled two-sample \(t\)-statistic under equal normal
variances.
\end{exercise}

\begin{exercise}[3]
Explain why paired data can have a much smaller standard error than an
independent-groups analysis.
\end{exercise}

\begin{exercise}[3]
Derive the two-proportion pooled standard error under
\(H_0:p_1=p_2\).
\end{exercise}

\begin{exercise}[3]
Derive the chi-square statistic for testing a normal variance.
\end{exercise}

\begin{exercise}[3]
Derive the \(F\)-statistic for comparing two normal population
variances.
\end{exercise}

\begin{exercise}[3]
Derive the expected-count formula

\[
E_{ij}
=
\frac{
n_{i\cdot}n_{\cdot j}
}{n}
\]

under independence.
\end{exercise}

\begin{exercise}[3]
Explain why a goodness-of-fit test with \(k\) fully specified categories
has \(k-1\) degrees of freedom.
\end{exercise}

\begin{exercise}[3]
Explain how estimating parameters from the same sample can reduce
goodness-of-fit degrees of freedom.
\end{exercise}

\begin{exercise}[3]
Derive the likelihood-ratio statistic

\[
-2\log\Lambda.
\]
\end{exercise}

\begin{exercise}[3]
Compare Wald, score, and likelihood-ratio tests conceptually.
\end{exercise}

\begin{exercise}[3]
State the Neyman--Pearson lemma and explain its optimality criterion.
\end{exercise}

\begin{exercise}[3]
Derive the power function for a one-sided known-variance normal mean
test.
\end{exercise}

\begin{exercise}[3]
Explain why power increases with sample size.
\end{exercise}

\begin{exercise}[3]
Explain why increasing \(\alpha\) increases both power and Type I error
risk.
\end{exercise}

\begin{exercise}[3]
Derive the sample-size planning formula for a one-sided \(z\)-test.
\end{exercise}

\begin{exercise}[3]
Explain why failure to reject \(H_0\) is not equivalent to evidence for
equivalence.
\end{exercise}

\begin{exercise}[3]
Explain the relationship between a two-sided hypothesis test and a
confidence interval.
\end{exercise}

\begin{exercise}[3]
Derive the Bonferroni family-wise error bound using the union bound.
\end{exercise}

\begin{exercise}[3]
Explain the difference between FWER and FDR.
\end{exercise}

\begin{exercise}[3]
Explain why permutation tests depend on exchangeability or the original
randomization design.
\end{exercise}

\begin{exercise}[3]
Explain how optional stopping can inflate Type I error.
\end{exercise}

\begin{exercise}[4]
Prove the Neyman--Pearson lemma.
\end{exercise}

\begin{exercise}[4]
Study monotone likelihood ratio families and derive UMP one-sided tests.
\end{exercise}

\begin{exercise}[4]
Investigate unbiased two-sided tests and UMP unbiased tests.
\end{exercise}

\begin{exercise}[4]
Prove Wilks' theorem under regularity conditions.
\end{exercise}

\begin{exercise}[4]
Study local asymptotic equivalence of Wald, score, and likelihood-ratio
tests.
\end{exercise}

\begin{exercise}[4]
Investigate exact binomial testing.
\end{exercise}

\begin{exercise}[4]
Study Fisher's exact test for \(2\times2\) tables.
\end{exercise}

\begin{exercise}[4]
Investigate permutation tests under unequal variances.
\end{exercise}

\begin{exercise}[4]
Study randomization inference under the potential-outcomes framework.
\end{exercise}

\begin{exercise}[4]
Investigate power and sample-size calculations for two-sample
\(t\)-tests.
\end{exercise}

\begin{exercise}[4]
Study equivalence and noninferiority testing.
\end{exercise}

\begin{exercise}[4]
Investigate sequential testing and alpha-spending.
\end{exercise}

\begin{exercise}[4]
Study the sequential probability ratio test.
\end{exercise}

\begin{exercise}[4]
Prove FWER control for Holm's procedure.
\end{exercise}

\begin{exercise}[4]
Study FDR control for the Benjamini--Hochberg procedure.
\end{exercise}

\begin{exercise}[4]
Investigate selective inference after model selection.
\end{exercise}

%%%%%%%%%%%%%%%%%%%%%%%%%%%%%%%%%%%%%%%%%%%%%%%%%%%%%%%%%%%%
\subsection{Conceptual Summary}
\label{subsec:hypothesis-testing-summary}
%%%%%%%%%%%%%%%%%%%%%%%%%%%%%%%%%%%%%%%%%%%%%%%%%%%%%%%%%%%%

A hypothesis test compares a null hypothesis

\[
\boxed{
H_0
}
\]

with an alternative hypothesis

\[
\boxed{
H_A.
}
\]

The significance level

\[
\boxed{
\alpha
}
\]

controls the probability of Type I error.

The probability of Type II error is

\[
\boxed{
\beta,
}
\]

and power is

\[
\boxed{
1-\beta.
}
\]

A typical standardized test statistic is

\[
\boxed{
T
=
\frac{
\text{estimate}-\text{null value}
}{
\text{null standard error}
}.
}
\]

The \(p\)-value is calculated under the null hypothesis and measures how
extreme the observed result is relative to the null sampling
distribution.

For a known-variance mean,

\[
\boxed{
Z
=
\frac{
\overline{X}-\mu_0
}{
\sigma/\sqrt{n}
}.
}
\]

For an unknown-variance normal mean,

\[
\boxed{
T
=
\frac{
\overline{X}-\mu_0
}{
S/\sqrt{n}
}
\sim
t_{n-1}.
}
\]

For one proportion,

\[
\boxed{
Z
=
\frac{
\widehat{p}-p_0
}{
\sqrt{
p_0(1-p_0)/n
}
}.
}
\]

For two independent means, Welch's statistic is

\[
\boxed{
T
=
\frac{
(\overline{X}_1-\overline{X}_2)-\Delta_0
}{
\sqrt{
S_1^2/n_1+S_2^2/n_2
}
}.
}
\]

For two proportions under equality,

\[
\boxed{
Z
=
\frac{
\widehat{p}_1-\widehat{p}_2
}{
\sqrt{
\widehat{p}_{pool}
(1-\widehat{p}_{pool})
(1/n_1+1/n_2)
}
}.
}
\]

For categorical count data,

\[
\boxed{
X^2
=
\sum
\frac{
(O-E)^2
}{
E
}.
}
\]

Hypothesis testing therefore converts a scientific claim into a
probability statement about how unusual the observed data would be under
a reference model.

%%%%%%%%%%%%%%%%%%%%%%%%%%%%%%%%%%%%%%%%%%%%%%%%%%%%%%%%%%%%
\subsection{Section Connections}
\label{subsec:hypothesis-testing-connections}
%%%%%%%%%%%%%%%%%%%%%%%%%%%%%%%%%%%%%%%%%%%%%%%%%%%%%%%%%%%%

\chapterconnection{
Hypothesis Testing uses the sampling distributions of
Section~\ref{sec:sampling-distributions}, the estimators of
Section~\ref{sec:statistical-estimation}, and the confidence intervals
of Section~\ref{sec:confidence-intervals}. Classical tests compare
observed estimators with null sampling distributions, while likelihood
theory produces Wald, score, and likelihood-ratio procedures. The next
section, Regression, extends these ideas to models relating response
variables to explanatory variables. Regression coefficients are
estimated from data, assigned standard errors, tested against null
values, and reported with confidence intervals, making regression a
direct continuation of the inference framework developed in Sections
8.3--8.6.
}

%%%%%%%%%%%%%%%%%%%%%%%%%%%%%%%%%%%%%%%%%%%%%%%%%%%%%%%%%%%%
% SECTION SOURCE TRACKING
%%%%%%%%%%%%%%%%%%%%%%%%%%%%%%%%%%%%%%%%%%%%%%%%%%%%%%%%%%%%

%------------------------------------------------------------
% 8.6 Hypothesis Testing
%------------------------------------------------------------
%
% Source basis:
%   - Standard mathematical statistics.
%   - Standard Neyman--Pearson testing theory.
%   - Standard classical and asymptotic hypothesis testing.
%
% Used for:
%   - null and alternative hypotheses;
%   - simple and composite hypotheses;
%   - one-sided and two-sided tests;
%   - test statistics;
%   - rejection regions;
%   - significance levels;
%   - Type I and Type II errors;
%   - power;
%   - p-values;
%   - statistical and practical significance;
%   - one-sample z-tests;
%   - one-sample t-tests;
%   - one-proportion tests;
%   - exact binomial tests;
%   - Welch two-sample tests;
%   - pooled two-sample tests;
%   - paired t-tests;
%   - two-proportion tests;
%   - normal variance tests;
%   - F-tests for ratios of variances;
%   - chi-square goodness-of-fit tests;
%   - chi-square tests of independence;
%   - tests of homogeneity;
%   - Pearson residuals;
%   - likelihood-ratio tests;
%   - Wilks theorem preview;
%   - Wald tests;
%   - score tests;
%   - Neyman--Pearson lemma;
%   - most powerful and UMP tests;
%   - power functions;
%   - sample-size planning;
%   - effect sizes;
%   - confidence-interval duality;
%   - equivalence and noninferiority previews;
%   - permutation and randomization tests;
%   - multiple testing;
%   - FWER;
%   - Bonferroni and Holm methods;
%   - FDR;
%   - Benjamini--Hochberg preview;
%   - optional stopping;
%   - sequential testing preview;
%   - applications and exercises.
%
% Notes:
%   - Regression is developed next in Section 8.7.
%   - Analysis of Variance is developed in Section 8.8.
%   - Experimental Design is developed in Section 8.9.
%   - Mathematical Statistics later develops Neyman--Pearson,
%     likelihood-ratio theory, and asymptotics more deeply.
%
%%%%%%%%%%%%%%%%%%%%%%%%%%%%%%%%%%%%%%%%%%%%%%%%%%%%%%%%%%%%